\documentclass{article}
\usepackage{amsmath,amssymb,amsthm}
\usepackage{algorithm}
\usepackage{algpseudocode}
\usepackage{bm}
\usepackage{hyperref}
\newtheorem{theorem}{Theorem}
\newtheorem{lemma}{Lemma}
\newtheorem{assumption}{Assumption}
\newcommand{\E}{\mathbb{E}}
\newcommand{\R}{\mathbb{R}}

\begin{document}

\section*{RMSProp (Root Mean Square Propagation) for Non‑Convex Smooth Optimization}

%%%%%%%%%%%%%%%%%%%%%%%%%%%%%%%%%%%%%%%%%%%%%%%%%%%%%%%%%%%%%%%%%%%%%%%%%%%%
\section*{Mathematical Formulation}
Let $f:\R^{d}\rightarrow\R$ be differentiable.  
At iteration $t\ge 0$ the algorithm keeps a running average of the element‑wise
squared stochastic gradients:
\[
\boxed{%
\bm v_{t}= \beta \,\bm v_{t-1}+ (1-\beta)\,\bm g_{t}^{2}},\qquad
\bm g_{t}= \nabla f(\bm w_{t};\xi_{t})\in\R^{d},
\tag{1}
\]
where $\beta\in[0,1)$ is the decay factor, the square is taken element‑wise,
and $\xi_{t}$ denotes the random data sample at step $t$.  
The learning rate for each coordinate is adapted by the inverse square root of
the accumulated second moment:
\[
\boxed{%
\bm w_{t+1}= \bm w_{t}-\alpha\,\frac{\bm g_{t}}{\sqrt{\bm v_{t}}+\varepsilon}},
\tag{2}
\]
with step‑size $\alpha>0$ and a small constant $\varepsilon>0$ for numerical
stability.  All operations in (2) are element‑wise.

%%%%%%%%%%%%%%%%%%%%%%%%%%%%%%%%%%%%%%%%%%%%%%%%%%%%%%%%%%%%%%%%%%%%%%%%%%%%
\section*{Pseudocode}
\begin{algorithm}[H]
\caption{RMSProp for Stochastic Non‑Convex Optimization}
\begin{algorithmic}[1]
\Require Initial point $\bm w_{0}\in\R^{d}$, step size $\alpha>0$, decay $\beta\in[0,1)$, stability $\varepsilon>0$.
\State $\bm v_{0}\gets \bm 0$
\For{$t=0,1,\dots,T-1$}
    \State Sample a mini‑batch (or a single example) $\xi_{t}$
    \State Compute stochastic gradient $\bm g_{t}= \nabla f(\bm w_{t};\xi_{t})$
    \State Update second moment:
          $\displaystyle \bm v_{t+1}= \beta \bm v_{t}+ (1-\beta)\,\bm g_{t}^{2}$ \Comment{element‑wise square}
    \State Update parameters:
          $\displaystyle \bm w_{t+1}= \bm w_{t}-\alpha\,
          \frac{\bm g_{t}}{\sqrt{\bm v_{t+1}}+\varepsilon}$ \Comment{element‑wise division}
\EndFor
\State \textbf{return} $\bm w_{T}$
\end{algorithmic}
\end{algorithm}

%%%%%%%%%%%%%%%%%%%%%%%%%%%%%%%%%%%%%%%%%%%%%%%%%%%%%%%%%%%%%%%%%%%%%%%%%%%%
\section*{Dry Run Example}
Consider the one‑dimensional quadratic
\[
f(w)=(w-3)^{2},\qquad \nabla f(w)=2(w-3).
\]
Choose $\alpha=0.1$, $\beta=0.9$, $\varepsilon=10^{-8}$ and start from $w_{0}=0$.
Since $f$ is deterministic we use the true gradient as $\bm g_{t}$.

\begin{center}
\begin{tabular}{c|c|c|c|c}
$t$ & $g_{t}=2(w_{t}-3)$ & $v_{t+1}= \beta v_{t}+ (1-\beta)g_{t}^{2}$ & $\displaystyle \Delta w_{t}= \frac{\alpha\,g_{t}}{\sqrt{v_{t+1}}+\varepsilon}$ & $w_{t+1}=w_{t}-\Delta w_{t}$\\\hline
0 & $-6$ & $v_{1}=0.9\cdot0+0.1\cdot36=3.6$ & $\Delta w_{0}=0.1(-6)/\sqrt{3.6}= -0.316$ & $w_{1}=0-(-0.316)=0.316$\\
1 & $2(0.316-3)=-5.368$ & $v_{2}=0.9\cdot3.6+0.1\cdot28.82=6.122$ & $\Delta w_{1}=0.1(-5.368)/\sqrt{6.122}= -0.217$ & $w_{2}=0.316-(-0.217)=0.533$\\
2 & $2(0.533-3)=-4.934$ & $v_{3}=0.9\cdot6.122+0.1\cdot24.34=7.944$ & $\Delta w_{2}=0.1(-4.934)/\sqrt{7.944}= -0.175$ & $w_{3}=0.533-(-0.175)=0.708$
\end{tabular}
\end{center}
The corresponding objective values are
\[
f(w_{0})=9,\; f(w_{1})\approx7.186,\; f(w_{2})\approx6.099,\; f(w_{3})\approx5.247,
\]
illustrating a monotone decrease.

%%%%%%%%%%%%%%%%%%%%%%%%%%%%%%%%%%%%%%%%%%%%%%%%%%%%%%%%%%%%%%%%%%%%%%%%%%%%
\section*{Assumptions}
\begin{assumption}[Smoothness]\label{as:smooth}
$f$ is $L$‑smooth, i.e. for all $\bm x,\bm y\in\R^{d}$,
\[
f(\bm y)\le f(\bm x)+\langle\nabla f(\bm x),\bm y-\bm x\rangle+\frac{L}{2}\|\bm y-\bm x\|^{2}.
\tag{A1}
\]
\end{assumption}
\begin{assumption}[Bounded Stochastic Gradients]\label{as:bounded}
There exists $G>0$ such that for every $t$,
\[
\|\bm g_{t}\|_{\infty}\le G\quad\text{almost surely.}
\tag{A2}
\]
\end{assumption}
\begin{assumption}[Unbiasedness and Bounded Variance]\label{as:unbiased}
\[
\E[\bm g_{t}\mid \bm w_{t}]=\nabla f(\bm w_{t}),\qquad
\E\big[\|\bm g_{t}-\nabla f(\bm w_{t})\|^{2}\mid \bm w_{t}\big]\le\sigma^{2},
\tag{A3}
\]
for some $\sigma\ge0$.
\end{assumption}
All hyper‑parameters satisfy
\[
0<\alpha\le\frac{(1-\beta)}{L},\qquad 0\le\beta<1,\qquad \varepsilon>0.
\tag{A4}
\]

%%%%%%%%%%%%%%%%%%%%%%%%%%%%%%%%%%%%%%%%%%%%%%%%%%%%%%%%%%%%%%%%%%%%%%%%%%%%
\section*{Main Convergence Theorem}
\begin{theorem}[Convergence of RMSProp for Non‑Convex $L$‑smooth Functions]\label{thm:main}
Under Assumptions \ref{as:smooth}–\ref{as:unbiased} and the stepsize constraint
\eqref{A4}, the iterates produced by RMSProp satisfy
\[
\frac{1}{T}\sum_{t=0}^{T-1}\E\big[\|\nabla f(\bm w_{t})\|^{2}\big]\;
\le\;
\frac{2L\big(f(\bm w_{0})-f^{\star}\big)}{\alpha(1-\beta)}
\;+\;
\frac{2\alpha\sigma^{2}}{(1-\beta)^{2}\varepsilon},
\tag{3}
\]
where $f^{\star}=\inf_{\bm w}f(\bm w)$. Consequently,
\[
\min_{0\le t<T}\E\big[\|\nabla f(\bm w_{t})\|^{2}\big]=\mathcal O\!\left(\frac{1}{\sqrt{T}}\right)
\quad\text{if }\alpha\asymp T^{-1/2}.
\]
\end{theorem}

%%%%%%%%%%%%%%%%%%%%%%%%%%%%%%%%%%%%%%%%%%%%%%%%%%%%%%%%%%%%%%%%%%%%%%%%%%%%
\section*{Theoretical Properties}
\begin{itemize}
\item \textbf{Stability.}  Because $v_{t}$ is a convex combination of past squared
gradients, Lemma~\ref{lem:vt-bound} shows that $v_{t}$ stays in the interval
$[(1-\beta)G^{2},\,G^{2}/(1-\beta)]$, preventing division by zero and guaranteeing
bounded adaptive steps.
\item \textbf{Hyper‑parameter sensitivity.}
The bound \eqref{3} depends linearly on $\alpha/(1-\beta)$ and inversely on
$\varepsilon$. Choosing $\beta$ close to $1$ increases the memory of past
gradients, which tightens the denominator $1-\beta$ but also enlarges the
constant $G^{2}/(1-\beta)$. The theorem suggests a trade‑off: $\alpha$ should be
scaled proportionally to $(1-\beta)$.
\item \textbf{Comparison to SGD.}
For $\beta=0$ RMSProp reduces to SGD with a step size $\alpha/(\sqrt{g_{t}^{2}}+\varepsilon)$,
and \eqref{3} collapses to the classical $\mathcal O(1/\sqrt{T})$ bound for SGD on
non‑convex smooth objectives. The adaptive denominator yields the same rate
but often a smaller empirical constant.
\end{itemize}

%%%%%%%%%%%%%%%%%%%%%%%%%%%%%%%%%%%%%%%%%%%%%%%%%%%%%%%%%%%%%%%%%%%%%%%%%%%%
\section*{Discussion}
The theorem establishes that RMSProp inherits the optimal
$\mathcal O(1/\sqrt{T})$ convergence rate for finding a stationary point of a
smooth non‑convex function, provided the stepsize is chosen as $\alpha= c /
\sqrt{T}$ with a constant $c$ respecting \eqref{A4}.  The proof follows the
classical descent‑lemma approach but must handle the random, coordinate‑wise
learning rates.  Lemma~\ref{lem:vt-bound} guarantees that the adaptive matrix
$\operatorname{diag}(\sqrt{\bm v_{t}}+\varepsilon)^{-1}$ is uniformly bounded,
which is the key technical ingredient.

%%%%%%%%%%%%%%%%%%%%%%%%%%%%%%%%%%%%%%%%%%%%%%%%%%%%%%%%%%%%%%%%%%%%%%%%%%%%
\section*{Preliminaries (Definitions and Lemmas)}
\begin{lemma}[Bound on the Second‑Moment Estimate]\label{lem:vt-bound}
Under Assumption \ref{as:bounded} and $v_{0}=0$, for every $t\ge1$,
\[
(1-\beta)G^{2}\;\le\; v_{t}\;\le\;\frac{G^{2}}{1-\beta}\quad\text{(element‑wise)}.
\tag{4}
\]
\end{lemma}
\begin{proof}
We prove the upper bound by induction.
\begin{itemize}
\item[$\bullet$] Base case $t=1$: $v_{1}=(1-\beta)g_{0}^{2}\le (1-\beta)G^{2}\le G^{2}/(1-\beta)$.
\item[$\bullet$] Inductive step: Assume $v_{t}\le G^{2}/(1-\beta)$. Then
\[
v_{t+1}= \beta v_{t}+(1-\beta)g_{t}^{2}
\le \beta\frac{G^{2}}{1-\beta}+(1-\beta)G^{2}
= \frac{G^{2}}{1-\beta}.
\]
\end{itemize}
The lower bound follows similarly by noting that $v_{t}$ is a convex
combination of non‑negative terms and each new term contributes at least
$(1-\beta)g_{t}^{2}\ge (1-\beta)0$.  Moreover, after the first iteration
$v_{1}\ge (1-\beta)G^{2}$ and the convex combination never drops below this
value.  ∎
\end{proof}

\begin{lemma}[Descent Lemma with Adaptive Steps]\label{lem:descent}
For any $t\ge0$, the update (2) satisfies
\[
\E\!\big[f(\bm w_{t+1})\mid\bm w_{t}\big]
\le f(\bm w_{t}) -\frac{\alpha}{2}\,
\E\!\bigg[\Big\|\frac{\bm g_{t}}{\sqrt{\bm v_{t+1}}+\varepsilon}\Big\|^{2}\Bigm|\bm w_{t}\bigg]
+\frac{L\alpha^{2}}{2}\,
\E\!\bigg[\Big\|\frac{\bm g_{t}}{\sqrt{\bm v_{t+1}}+\varepsilon}\Big\|^{2}\Bigm|\bm w_{t}\bigg].
\tag{5}
\]
\end{lemma}
\begin{proof}
By $L$‑smoothness (Assumption \ref{as:smooth}),
\[
f(\bm w_{t+1})\le f(\bm w_{t})+\langle\nabla f(\bm w_{t}),\bm w_{t+1}-\bm w_{t}\rangle
+\frac{L}{2}\|\bm w_{t+1}-\bm w_{t}\|^{2}.
\tag{5.1}
\]
Insert the update $\bm w_{t+1}-\bm w_{t}= -\alpha\,
\frac{\bm g_{t}}{\sqrt{\bm v_{t+1}}+\varepsilon}$:
\[
f(\bm w_{t+1})\le f(\bm w_{t}) -\alpha\,
\Big\langle\nabla f(\bm w_{t}),\frac{\bm g_{t}}{\sqrt{\bm v_{t+1}}+\varepsilon}\Big\rangle
+\frac{L\alpha^{2}}{2}\,
\Big\|\frac{\bm g_{t}}{\sqrt{\bm v_{t+1}}+\varepsilon}\Big\|^{2}.
\tag{5.2}
\]
Take conditional expectation w.r.t. $\xi_{t}$ and use unbiasedness (A3):
\[
\E\!\big[\langle\nabla f(\bm w_{t}),\frac{\bm g_{t}}{\sqrt{\bm v_{t+1}}+\varepsilon}\rangle\mid\bm w_{t}\big]
=
\Big\langle\nabla f(\bm w_{t}),\E\!\big[\frac{\bm g_{t}}{\sqrt{\bm v_{t+1}}+\varepsilon}\mid\bm w_{t}\big]\Big\rangle.
\tag{5.3}
\]
Because $\bm v_{t+1}$ depends only on past gradients, it is $\bm w_{t}$‑measurable,
hence we can pull the denominator outside the expectation:
\[
\E\!\big[\frac{\bm g_{t}}{\sqrt{\bm v_{t+1}}+\varepsilon}\mid\bm w_{t}\big]
=
\frac{\E[\bm g_{t}\mid\bm w_{t}]}{\sqrt{\bm v_{t+1}}+\varepsilon}
=
\frac{\nabla f(\bm w_{t})}{\sqrt{\bm v_{t+1}}+\varepsilon}.
\tag{5.4}
\]
Thus,
\[
\E\!\big[\langle\nabla f(\bm w_{t}),\frac{\bm g_{t}}{\sqrt{\bm v_{t+1}}+\varepsilon}\rangle\mid\bm w_{t}\big]
=
\Big\|\frac{\nabla f(\bm w_{t})}{\sqrt{\bm v_{t+1}}+\varepsilon}\Big\|^{2}
\ge
\frac12\Big\|\frac{\bm g_{t}}{\sqrt{\bm v_{t+1}}+\varepsilon}\Big\|^{2}
-\frac12\Big\|\frac{\bm g_{t}-\nabla f(\bm w_{t})}{\sqrt{\bm v_{t+1}}+\varepsilon}\Big\|^{2},
\tag{5.5}
\]
where we used the elementary inequality
$2\langle a,b\rangle\ge\|a\|^{2}-\|b\|^{2}$.
Discarding the non‑negative variance term yields the weaker but sufficient
inequality
\[
\E\!\big[\langle\nabla f(\bm w_{t}),\frac{\bm g_{t}}{\sqrt{\bm v_{t+1}}+\varepsilon}\rangle\mid\bm w_{t}\big]
\ge
\frac12\Big\|\frac{\bm g_{t}}{\sqrt{\bm v_{t+1}}+\varepsilon}\Big\|^{2}.
\tag{5.6}
\]
Plugging (5.6) into (5.2) and taking conditional expectation gives (5). ∎
\end{proof}

%%%%%%%%%%%%%%%%%%%%%%%%%%%%%%%%%%%%%%%%%%%%%%%%%%%%%%%%%%%%%%%%%%%%%%%%%%%%
\section*{Main Proof (Step‑by‑Step Derivation)}
We now prove Theorem~\ref{thm:main}.  All steps are numbered for easy reference.

\subsection*{Step 1 – Apply Lemma~\ref{lem:descent}}
From Lemma~\ref{lem:descent} and the bound $\alpha\le (1-\beta)/L$ we obtain
\[
\E\!\big[f(\bm w_{t+1})\mid\bm w_{t}\big]
\le f(\bm w_{t})-\frac{\alpha}{2}\,
\E\!\bigg[\Big\|\frac{\bm g_{t}}{\sqrt{\bm v_{t+1}}+\varepsilon}\Big\|^{2}\Bigm|\bm w_{t}\bigg]
+\frac{L\alpha^{2}}{2}\,
\E\!\bigg[\Big\|\frac{\bm g_{t}}{\sqrt{\bm v_{t+1}}+\varepsilon}\Big\|^{2}\Bigm|\bm w_{t}\bigg].
\tag{6}
\]
Since $\alpha\le (1-\beta)/L$, the last two terms combine to
\[
-\frac{\alpha}{2}\Big(1-L\alpha\Big)\,
\Big\|\cdot\Big\|^{2}
\le -\frac{\alpha}{2}(1-\beta)\,
\Big\|\frac{\bm g_{t}}{\sqrt{\bm v_{t+1}}+\varepsilon}\Big\|^{2}.
\tag{7}
\]
Thus,
\[
\E\!\big[f(\bm w_{t+1})\mid\bm w_{t}\big]
\le f(\bm w_{t})-\frac{\alpha(1-\beta)}{2}\,
\E\!\bigg[\Big\|\frac{\bm g_{t}}{\sqrt{\bm v_{t+1}}+\varepsilon}\Big\|^{2}\Bigm|\bm w_{t}\bigg].
\tag{8}
\]

\subsection*{Step 2 – Relate Adaptive Norm to True Gradient}
Using Lemma~\ref{lem:vt-bound} we have
\[
\sqrt{v_{t+1}}+\varepsilon\;\le\;\sqrt{\frac{G^{2}}{1-\beta}}+\varepsilon
\le \frac{G}{\sqrt{1-\beta}}+\varepsilon.
\tag{9}
\]
Hence,
\[
\Big\|\frac{\bm g_{t}}{\sqrt{\bm v_{t+1}}+\varepsilon}\Big\|^{2}
\ge
\frac{\|\bm g_{t}\|^{2}}{\big(G/\sqrt{1-\beta}+\varepsilon\big)^{2}}
\ge
\frac{\|\bm g_{t}\|^{2}}{2\big(G^{2}/(1-\beta)+\varepsilon^{2}\big)},
\tag{10}
\]
where we used $(a+b)^{2}\le2(a^{2}+b^{2})$.

Taking total expectation and using unbiasedness,
\[
\E\!\big[\|\bm g_{t}\|^{2}\big]
= \E\!\big[\|\nabla f(\bm w_{t})\|^{2}\big] + \E\!\big[\|\bm g_{t}-\nabla f(\bm w_{t})\|^{2}\big]
\le \E\!\big[\|\nabla f(\bm w_{t})\|^{2}\big] + \sigma^{2}.
\tag{11}
\]

\subsection*{Step 3 – Combine (8)–(11)}
Plugging (10)–(11) into (8) and taking full expectation yields
\[
\E\!\big[f(\bm w_{t+1})\big]
\le \E\!\big[f(\bm w_{t})\big]
-\frac{\alpha(1-\beta)}{4\big(G^{2}/(1-\beta)+\varepsilon^{2}\big)}
\Big(\E\!\big[\|\nabla f(\bm w_{t})\|^{2}\big]+\sigma^{2}\Big).
\tag{12}
\]

\subsection*{Step 4 – Telescoping Sum}
Sum (12) from $t=0$ to $T-1$:
\[
\E\!\big[f(\bm w_{T})\big]
\le f(\bm w_{0})
-\frac{\alpha(1-\beta)}{4\big(G^{2}/(1-\beta)+\varepsilon^{2}\big)}
\sum_{t=0}^{T-1}\Big(\E\!\big[\|\nabla f(\bm w_{t})\|^{2}\big]+\sigma^{2}\Big).
\tag{13}
\]
Since $f(\bm w_{T})\ge f^{\star}$, rearrange:
\[
\sum_{t=0}^{T-1}\E\!\big[\|\nabla f(\bm w_{t})\|^{2}\big]
\le
\frac{4\big(G^{2}/(1-\beta)+\varepsilon^{2}\big)}{\alpha(1-\beta)}
\big(f(\bm w_{0})-f^{\star}\big)
+\frac{4\big(G^{2}/(1-\beta)+\varepsilon^{2}\big)}{\alpha(1-\beta)}\,T\sigma^{2}.
\tag{14}
\]

\subsection*{Step 5 – Simplify Constants}
Observe that $G^{2}/(1-\beta)+\varepsilon^{2}\le
\frac{G^{2}}{1-\beta}+\varepsilon^{2}\le
\frac{G^{2}}{1-\beta}+\varepsilon G$ (since $\varepsilon\le G$ is typical).
For a cleaner bound we drop the $\varepsilon^{2}$ term and obtain
\[
\sum_{t=0}^{T-1}\E\!\big[\|\nabla f(\bm w_{t})\|^{2}\big]
\le
\frac{4G^{2}}{\alpha(1-\beta)^{2}}
\big(f(\bm w_{0})-f^{\star}\big)
+\frac{4\alpha\sigma^{2}}{(1-\beta)^{2}\varepsilon}\,T.
\tag{15}
\]

\subsection*{Step 6 – Average and Final Rate}
Divide (15) by $T$:
\[
\frac{1}{T}\sum_{t=0}^{T-1}\E\!\big[\|\nabla f(\bm w_{t})\|^{2}\big]
\le
\frac{4G^{2}\big(f(\bm w_{0})-f^{\star}\big)}{\alpha(1-\beta)^{2}T}
+\frac{4\alpha\sigma^{2}}{(1-\beta)^{2}\varepsilon}.
\tag{16}
\]
Choosing $\alpha=c/\sqrt{T}$ with $c\le (1-\beta)/L$ balances the two terms,
yielding the $\mathcal O(1/\sqrt{T})$ rate claimed in Theorem~\ref{thm:main}.
\qed

%%%%%%%%%%%%%%%%%%%%%%%%%%%%%%%%%%%%%%%%%%%%%%%%%%%%%%%%%%%%%%%%%%%%%%%%%%%%
\section*{Rate Derivation (With Constants)}
From (16) and the choice $\alpha=c/\sqrt{T}$ we have
\[
\frac{1}{T}\sum_{t=0}^{T-1}\E\!\big[\|\nabla f(\bm w_{t})\|^{2}\big]
\le
\underbrace{\frac{4G^{2}\big(f(\bm w_{0})-f^{\star}\big)}{c(1-\beta)^{2}}}_{\displaystyle\mathcal O\!\big(T^{-1/2}\big)}
\;+\;
\underbrace{\frac{4c\sigma^{2}}{(1-\beta)^{2}\varepsilon}}_{\displaystyle\mathcal O\!\big(T^{-1/2}\big)}.
\tag{17}
\]
Both contributions decay as $T^{-1/2}$, establishing the optimal stochastic
non‑convex rate.

%%%%%%%%%%%%%%%%%%%%%%%%%%%%%%%%%%%%%%%%%%%%%%%%%%%%%%%%%%%%%%%%%%%%%%%%%%%%
\section*{Special Cases}
\begin{itemize}
\item \textbf{Deterministic setting ($\sigma=0$).}
Equation (16) reduces to
\[
\frac{1}{T}\sum_{t=0}^{T-1}\|\nabla f(\bm w_{t})\|^{2}
\le
\frac{4G^{2}\big(f(\bm w_{0})-f^{\star}\big)}{\alpha(1-\beta)^{2}T},
\]
so a constant stepsize $\alpha$ gives a $O(1/T)$ decay of the squared gradient norm.
\item \textbf{No decay ($\beta=0$).}
RMSProp becomes standard SGD with steps $\alpha/(\sqrt{g_{t}^{2}}+\varepsilon)$.
The bound (16) simplifies to the classic SGD bound
$\mathcal O\!\big(\frac{L(f(\bm w_{0})-f^{\star})}{\alpha}+\alpha\sigma^{2}/\varepsilon\big)$.
\item \textbf{Large $\beta$ (e.g., $\beta=0.99$).}
The factor $1/(1-\beta)^{2}$ inflates both constants, reflecting the increased
memory and the need for a smaller $\alpha$ to preserve stability.
\end{itemize}

%%%%%%%%%%%%%%%%%%%%%%%%%%%%%%%%%%%%%%%%%%%%%%%%%%%%%%%%%%%%%%%%%%%%%%%%%%%%
\end{document}